%%%%%%%%%%%%%%%%
%%% gen 38 %%%
%%%%%%%%%%%%%%%%

\Chapter{38}

\Verse{1}
{
	\hJ{וַיְהִי בָּעֵת הַהִוא וַיֵּרֶד יְהוּדָה מֵאֵת אֶחָיו וַיֵּט עַד אִישׁ עֲדֻלָּמִי וּשְׁמוֹ חִירָה׃}
}
{
	\eJ{
		And it was at that time, and Judah went down from
		his brothers and turned to an Adullamite man, and his name
		was Hirah.
	}
}
{

%	NOTES TO PREPARE FOR THE SCENE:
%
%	\begin{quote}
%	A breach is a break, a tear, or a failure to obey a law, rule, promise, or duty.
%	\end{quote}
%
%	Onan — אונן → 𐤀𐤅𐤍𐤍
%	(Sex with Tamar, violation, killed for it.)
%
%	Amnon — אמנון → 𐤀𐤌𐤍𐤅𐤍
%	(Sex with Tamar, violation, killed for it.)
%
%	Shelah — שלה → 𐤔𐤋𐤄
%	(Supposed to marry Tamar, doesn't.)
%
%	Solomon — שלמה → 𐤔𐤋𐤌𐤄
%	(King around the time this is written.)
%
%	Perez — פרץ → 𐤐𐤓𐤑
%	(Twin w/o red cord, pushes ahead of brother.)
%
%	Zerah — זרח → 𐤆𐤓𐤇
%	(Twin w/ red cord, second fully out.)
%
%	Breach — פרץ → 𐤐𐤓𐤑
%
%	Seed — זרע → 𐤆𐤓𐤏
%
%	Tamar — תמר → 𐤕𐤌𐤓 (TMR)
%
%	She said — תאמר → 𐤕𐤀𐤌𐤓 (T’MR)
%
%	Maacah — מעכה → 𐤌𐤏𐤊𐤄
%	(Daughter of Absalom, married Rehoboam.)
%
%	YHWH — יהוה → 𐤉𐤄𐤅𐤄
%
%	David — דוד → 𐤃𐤅𐤃
%
%	Judah — יהודה → 𐤉𐤄𐤅𐤃𐤄
%
%	Bathsheba — בתשבע → 𐤁𐤕𐤔𐤁𐤏
%
%	Daughter of Shua — בתשוע → 𐤁𐤕𐤔𐤅𐤏
%
%	The only difference between Bathsheba and Daughter of Shua is writing a letter that looks like this 𐤅 in Paleo instead of a letter that looks like this 𐡁 in Aramaic.
%
%	When the Torah is redacted (following the account in 2 Esdras) it would have been transliterated from Paleo Hebrew to Aramaic. According to 2 Esdras, the scribes copied it quickly, using ``letters they did not know.''
%
%	\begingroup
%	\scriptsize
%	\setlength{\tabcolsep}{0.35em}
%	\Table{@{}p{0.13\linewidth}p{0.18\linewidth}p{0.20\linewidth}p{0.13\linewidth}@{}}{
%	\hline
%		\textbf{Letter} & \textbf{Paleo} & \textbf{Aramaic} & \textbf{Hebrew} \\
%	\hline
%		ʾ & 𐤀 & 𐡀 & \heb{א} \\
%	\hline
%		B & 𐤁 & 𐡁 & \heb{ב} \\
%	\hline
%		G & 𐤂 & 𐡂 & \heb{ג} \\
%	\hline
%		D & 𐤃 & 𐡃 & \heb{ד} \\
%	\hline
%		H & 𐤄 & 𐡄 & \heb{ה} \\
%	\hline
%		W & 𐤅 & 𐡅 & \heb{ו} \\
%	\hline
%		Z & 𐤆 & 𐡆 & \heb{ז} \\
%	\hline
%		Ḥ & 𐤇 & 𐡇 & \heb{ח} \\
%	\hline
%		Ṭ & 𐤈 & 𐡈 & \heb{ט} \\
%	\hline
%		Y & 𐤉 & 𐡉 & \heb{י} \\
%	\hline
%		K & 𐤊 & 𐡊 & \heb{כ} \\
%	\hline
%		L & 𐤋 & 𐡋 & \heb{ל} \\
%	\hline
%		M & 𐤌 & 𐡌 & \heb{מ} \\
%	\hline
%		N & 𐤍 & 𐡍 & \heb{נ} \\
%	\hline
%		S & 𐤎 & 𐡎 & \heb{ס} \\
%	\hline
%		ʿ & 𐤏 & 𐡏 & \heb{ע} \\
%	\hline
%		P & 𐤐 & 𐡐 & \heb{פ} \\
%	\hline
%		Ṣ & 𐤑 & 𐡑 & \heb{צ} \\
%	\hline
%		Q & 𐤒 & 𐡒 & \heb{ק} \\
%	\hline
%		R & 𐤓 & 𐡓 & \heb{ר} \\
%	\hline
%		Š & 𐤔 & 𐡔 & \heb{ש} \\
%	\hline
%		T & 𐤕 & 𐡕 & \heb{ת} \\
%	\hline
%	}
%	\endgroup

%	\footnote{
%		Judah. The eponymous ancestor of the Jews, Judah is the most prominent
%		of the brothers in the Joseph stories, and here he is the only one of
%		Jacob’s sons besides Joseph to have a separate story about him. Some say
%		that the word means “to be thankful.” Its derivation is unknown.
%		Footnote 38:1. Judah went down. The next chapter begins with “Joseph had
%		been brought down.” The contrast is blatant (even more so in the Hebrew)
%		between Judah’s independence and Joseph’s weakness: Judah went down, and
%		Joseph was brought down.
%	}
}

\Verse{2}
{
	\hJ{וַיַּרְא שָׁם יְהוּדָה בַּת אִישׁ כְּנַעֲנִי וּשְׁמוֹ שׁוּעַ וַיִּקָּחֶהָ וַיָּבֹא אֵלֶיהָ׃}
}
{
	\eJ{
		And Judah saw a daughter of a Canaanite man there, and his
		name was Shua. And he took her and came to her.
	}
}
{
	\Table{|p{0.16\linewidth}|p{0.25\linewidth}|p{0.49\linewidth}|}{
		\hline
		\textbf{Proper name} & \textbf{Hebrew} & \textbf{Possible Meaning} \\
		\hline
		Judah & \heb{יְהוּדָה} & See v.1. \\
		\hline
		Shua & \heb{שׁוּעַ} (\emph{Šûaʿ})
			& Possibly related to \heb{שוע}, “cry for help,” but uncertain. \\
		\hline
	}
}

\Verse{3}
{
	\hJ{וַתַּהַר וַתֵּלֶד בֵּן וַיִּקְרָא אֶת שְׁמוֹ עֵר׃}
}
{
	\eJ{
		And she became pregnant and gave birth to a son, and he
		called his name Er.
	}
}
{
	\Table{|p{0.16\linewidth}|p{0.25\linewidth}|p{0.49\linewidth}|}{
		\hline
		\textbf{Proper name} & \textbf{Hebrew} & \textbf{Possible Meaning} \\
		\hline
		Er & \heb{עֵר} (\emph{ʿĒr})
			& Traditionally associated with \heb{עור}, “be awake,” hence “watchful.” \\
		\hline
	}
}

\Verse{4}
{
	\hJ{וַתַּהַר עוֹד וַתֵּלֶד בֵּן וַתִּקְרָא אֶת שְׁמוֹ אוֹנָן׃}
}
{
	\eJ{
		And she became pregnant again and gave birth to a son, and
		she called his name Onan.
	}
}
{
	\Table{|p{0.16\linewidth}|p{0.25\linewidth}|p{0.49\linewidth}|}{
		\hline
		\textbf{Proper name} & \textbf{Hebrew} & \textbf{Possible Meaning} \\
		\hline
		Onan & \heb{אוֹנָן} (\emph{ʾÔnān})
			& Related to \heb{אוֹן}, “strength/vigor.” \\
		\hline
	}
}

\Verse{5}
{
	\hJ{וַתֹּסֶף עוֹד וַתֵּלֶד בֵּן וַתִּקְרָא אֶת שְׁמוֹ שֵׁלָה וְהָיָה בִכְזִיב בְּלִדְתָּהּ אֹתוֹ׃}
}
{
	\eJ{
		And she proceeded again to give birth to a son, and she
		called his name Shelah. (And he was at Chezib when she
		gave birth to him.)
	}
}
{
	\Table{|p{0.16\linewidth}|p{0.25\linewidth}|p{0.49\linewidth}|}{
		\hline
		\textbf{Proper name} & \textbf{Hebrew} & \textbf{Possible Meaning} \\
		\hline
		Shelah & \heb{שֵׁלָה} (\emph{Šēlāh})
			& Possibly “petition/request”; the derivation is uncertain. \\
		\hline
	}
}

\Verse{6}
{
	\hJ{וַיִּקַּח יְהוּדָה אִשָּׁה לְעֵר בְּכוֹרוֹ וּשְׁמָהּ תָּמָר׃}
}
{
	\eJ{
		And Judah took a wife for Er, his firstborn, and her name
		was Tamar.
	}
}
{
	\Table{|p{0.16\linewidth}|p{0.25\linewidth}|p{0.49\linewidth}|}{
		\hline
		\textbf{Proper name} & \textbf{Hebrew} & \textbf{Possible Meaning} \\
		\hline
		Judah & \heb{יְהוּדָה} & See v.1. \\
		\hline
		Er & \heb{עֵר} & See v.3. \\
		\hline
		Tamar & \heb{תָּמָר} (\emph{Tāmār}) & “Date palm.” \\
		\hline
	}
}

\Verse{7}
{
	\hJ{וַיְהִי עֵר בְּכוֹר יְהוּדָה רַע בְּעֵינֵי יְהֹוָה וַיְמִתֵהוּ יְהֹוָה׃}
}
{
	\eJ{
		And Er, Judah’s firstborn, was bad in YHWH’s eyes, and
		YHWH killed him.
	}
}
{
	\Table{|p{0.16\linewidth}|p{0.25\linewidth}|p{0.49\linewidth}|}{
		\hline
		\textbf{Proper name} & \textbf{Hebrew} & \textbf{Possible Meaning} \\
		\hline
		Er & \heb{עֵר} & See v.3. \\
		\hline
		Judah & \heb{יְהוּדָה} & See v.1. \\
		\hline
		YHWH & \heb{יְהֹוָה} (\emph{YHWH})
			& The divine name, associated in Exodus 3:14 with \heb{היה}, “be.” \\
		\hline
	}
}

\Verse{8}
{
	\hJ{וַיֹּאמֶר יְהוּדָה לְאוֹנָן בֹּא אֶל אֵשֶׁת אָחִיךָ וְיַבֵּם אֹתָהּ וְהָקֵם זֶרַע לְאָחִיךָ׃}
}
{
	\eJ{
		And Judah said to Onan, “Come to your brother’s wife and
		couple as a brother-in-law with her and raise seed for
		your brother.”
	}
}
{
	\Table{|p{0.16\linewidth}|p{0.25\linewidth}|p{0.49\linewidth}|}{
		\hline
		\textbf{Proper name} & \textbf{Hebrew} & \textbf{Possible Meaning} \\
		\hline
		Judah & \heb{יְהוּדָה} & See v.1. \\
		\hline
		Onan & \heb{אוֹנָן} & See v.4. \\
		\hline
	}
%	\footnote{
%		raise seed for your brother. This is known as the law of levirate
%		marriage, from the Latin levir, meaning brother-in-law. See Deuteronomy
%		25 and the comments on it.
%	}
}

\Verse{9}
{
	\hJ{וַיֵּדַע אוֹנָן כִּי לֹּא לוֹ יִהְיֶה הַזָּרַע וְהָיָה אִם בָּא אֶל אֵשֶׁת אָחִיו וְשִׁחֵת אַרְצָה לְבִלְתִּי נְתן זֶרַע לְאָחִיו׃}
}
{
	\eJ{
		And Onan knew that the seed would not be his. And it was
		when he came to his brother’s wife: and he spent on the
		ground so as not to give seed for his brother.
	}
}
{
	\Table{|p{0.16\linewidth}|p{0.25\linewidth}|p{0.49\linewidth}|}{
		\hline
		\textbf{Proper name} & \textbf{Hebrew} & \textbf{Possible Meaning} \\
		\hline
		Onan & \heb{אוֹנָן} & See v.4. \\
		\hline
	}
}

\Verse{10}
{
	\hJ{וַיֵּרַע בְּעֵינֵי יְהֹוָה אֲשֶׁר עָשָׂה וַיָּמֶת גַּם אֹתוֹ׃}
}
{
	\eJ{
		And what he did was bad in YHWH’s eyes, and He killed him,
		too.
	}
}
{
	\Table{|p{0.16\linewidth}|p{0.25\linewidth}|p{0.49\linewidth}|}{
		\hline
		\textbf{Proper name} & \textbf{Hebrew} & \textbf{Possible Meaning} \\
		\hline
		YHWH & \heb{יְהֹוָה} & See v.7. \\
		\hline
	}
}

\Verse{11}
{
	\hJ{וַיֹּאמֶר יְהוּדָה לְתָמָר כַּלָּתוֹ שְׁבִי אַלְמָנָה בֵית אָבִיךְ עַד יִגְדַּל שֵׁלָה בְנִי כִּי אָמַר פֶּן יָמוּת גַּם הוּא כְּאֶחָיו וַתֵּלֶךְ תָּמָר וַתֵּשֶׁב בֵּית אָבִיהָ׃}
}
{
	\eJ{
		And Judah said to Tamar, his daughter-in-law, “Live as a
		widow at your father’s house until my son Shelah grows up”
		(because he said, “Or else he, too, will die like his
		brothers”). And Tamar went and lived at her father’s
		house.
	}
}
{
	\Table{|p{0.16\linewidth}|p{0.25\linewidth}|p{0.49\linewidth}|}{
		\hline
		\textbf{Proper name} & \textbf{Hebrew} & \textbf{Possible Meaning} \\
		\hline
		Judah & \heb{יְהוּדָה} & See v.1. \\
		\hline
		Tamar & \heb{תָּמָר} & See v.6. \\
		\hline
		Shelah & \heb{שֵׁלָה} & See v.5. \\
		\hline
	}
}

\Verse{12}
{
	\hJ{וַיִּרְבּוּ הַיָּמִים וַתָּמת בַּת שׁוּעַ אֵשֶׁת יְהוּדָה וַיִּנָּחֶם יְהוּדָה וַיַּעַל עַל גֹּזְזֵי צֹאנוֹ הוּא וְחִירָה רֵעֵהוּ הָעֲדֻלָּמִי תִּמְנָתָה׃}
}
{
	\eJ{
		And the days were many, and Judah’s wife, the daughter of
		Shua, died. And Judah was consoled. And he went up to his
		sheepshearers, he and his friend Hirah, the Adullamite, to
		Timnah.
	}
}
{
	\Table{|p{0.16\linewidth}|p{0.25\linewidth}|p{0.49\linewidth}|}{
		\hline
		\textbf{Proper name} & \textbf{Hebrew} & \textbf{Possible Meaning} \\
		\hline
		Shua & \heb{שׁוּעַ} & See v.2. \\
		\hline
		Judah & \heb{יְהוּדָה} & See v.1. \\
		\hline
		Hirah & \heb{חִירָה} & See v.1. \\
		\hline
	}
}

\Verse{13}
{
	\hJ{וַיֻּגַּד לְתָמָר לֵאמֹר הִנֵּה חָמִיךְ עֹלֶה תִמְנָתָה לָגֹז צֹאנוֹ׃}
}
{
	\eJ{
		And it was told to Tamar, saying, “Here, your
		father-in-law is going up to Timnah to shear his sheep.”
	}
}
{
	\Table{|p{0.16\linewidth}|p{0.25\linewidth}|p{0.49\linewidth}|}{
		\hline
		\textbf{Proper name} & \textbf{Hebrew} & \textbf{Possible Meaning} \\
		\hline
		Tamar & \heb{תָּמָר} & See v.6. \\
		\hline
	}
}

\Verse{14}
{
	\hJ{וַתָּסַר בִּגְדֵי אַלְמְנוּתָהּ מֵעָלֶיהָ וַתְּכַס בַּצָּעִיף וַתִּתְעַלָּף וַתֵּשֶׁב בְּפֶתַח עֵינַיִם אֲשֶׁר עַל דֶּרֶךְ תִּמְנָתָה כִּי רָאֲתָה כִּי גָדַל שֵׁלָה וְהִוא לֹא נִתְּנָה לוֹ לְאִשָּׁה׃}
}
{
	\eJ{
		And she took off her widowhood clothes from on her and
		covered herself with a veil and wrapped herself and sat in
		a visible place that was on the road to Timnah, because
		she saw that Shelah had grown up and she had not been
		given to him as a wife.
	}
}
{
	\Table{|p{0.16\linewidth}|p{0.25\linewidth}|p{0.49\linewidth}|}{
		\hline
		\textbf{Proper name} & \textbf{Hebrew} & \textbf{Possible Meaning} \\
		\hline
		Shelah & \heb{שֵׁלָה} & See v.5. \\
		\hline
	}
}

\Verse{15}
{
	\hJ{וַיִּרְאֶהָ יְהוּדָה וַיַּחְשְׁבֶהָ לְזוֹנָה כִּי כִסְּתָה פָּנֶיהָ׃}
}
{
	\eJ{
		And Judah saw her and thought her to be a prostitute
		because she had covered her face.
	}
}
{
}

\Verse{16}
{
	\hJ{וַיֵּט אֵלֶיהָ אֶל הַדֶּרֶךְ וַיֹּאמֶר הָבָה נָּא אָבוֹא אֵלַיִךְ כִּי לֹא יָדַע כִּי כַלָּתוֹ הִוא וַתֹּאמֶר מַה תִּתֶּן לִי כִּי תָבוֹא אֵלָי׃}
}
{
	\eJ{
		And he turned to her by the road. And he said, “Come on.
		Let me come to you,” because he did not know that she was
		his daughter-in-law. And she said, “What will you give me
		when you come to me?”
	}
}
{
}

\Verse{17}
{
	\hJ{וַיֹּאמֶר אָנֹכִי אֲשַׁלַּח גְּדִי עִזִּים מִן הַצֹּאן וַתֹּאמֶר אִם תִּתֵּן עֵרָבוֹן עַד שׁלְחֶךָ׃}
}
{
	\eJ{
		And he said, “I’ll have a goat kid sent from the flock.”
		And she said, “If you’ll give a pledge until you send it.”
	}
}
{
}

\Verse{18}
{
	\hJ{וַיֹּאמֶר מָה הָעֵרָבוֹן אֲשֶׁר אֶתֶּן לָךְ וַתֹּאמֶר חֹתָמְךָ וּפְתִילֶךָ וּמַטְּךָ אֲשֶׁר בְּיָדֶךָ וַיִּתֶּן לָהּ וַיָּבֹא אֵלֶיהָ וַתַּהַר לוֹ׃}
}
{
	\eJ{
		And he said, “What is the pledge that I’ll give you?” And
		she said, “Your seal and your cord and your staff that’s
		in your hand.” And he gave them to her, and he came to
		her, and she became pregnant by him.
	}
}
{
	In J, women get power by hook or by crook.\fB{Lit. ``by hooker by crook'' in the original text.}
}

\Verse{19}
{
	\hJ{וַתָּקם וַתֵּלֶךְ וַתָּסַר צְעִיפָהּ מֵעָלֶיהָ וַתִּלְבַּשׁ בִּגְדֵי אַלְמְנוּתָהּ׃}
}
{
	\eJ{
		And she got up and went and took off her veil from on her
		and put on her widowhood clothes.
	}
}
{
}

\Verse{20}
{
	\hJ{וַיִּשְׁלַח יְהוּדָה אֶת גְּדִי הָעִזִּים בְּיַד רֵעֵהוּ הָעֲדֻלָּמִי לָקַחַת הָעֵרָבוֹן מִיַּד הָאִשָּׁה וְלֹא מְצָאָהּ׃}
}
{
	\eJ{
		And Judah sent the goat kid by the hand of his friend the
		Adullamite to take back the pledge from the woman’s hand,
		and he did not find her.
	}
}
{
	\Table{|p{0.16\linewidth}|p{0.25\linewidth}|p{0.49\linewidth}|}{
		\hline
		\textbf{Proper name} & \textbf{Hebrew} & \textbf{Possible Meaning} \\
		\hline
		Judah & \heb{יְהוּדָה} & See v.1. \\
		\hline
	}
}

\Verse{21}
{
	\hJ{וַיִּשְׁאַל אֶת אַנְשֵׁי מְקֹמָהּ לֵאמֹר אַיֵּה הַקְּדֵשָׁה הִוא בָעֵינַיִם עַל הַדָּרֶךְ וַיֹּאמְרוּ לֹא הָיְתָה בָזֶה קְדֵשָׁה׃}
}
{
	\eJ{
		And he asked the people of her place, saying, “Where is
		the sacred prostitute? She was visibly by the road.” And
		they said, “There was no sacred prostitute here.”
	}
}
{
%	\footnote{
%		sacred prostitute. On sacred prostitution, see the comment on Deut
%		23:18. The immediate point in this episode is that Judah is described as
%		thinking that Tamar is a prostitute (znh) (38:15), but now his friend
%		uses what is seemingly a more refined word, “sacred prostitute” (qdšh)—
%		which may also denote a higher level of prostitute—when he discreetly
%		inquires about her.
%	}
}

\Verse{22}
{
	\hJ{וַיָּשׁב אֶל יְהוּדָה וַיֹּאמֶר לֹא מְצָאתִיהָ וְגַם אַנְשֵׁי הַמָּקוֹם אָמְרוּ לֹא הָיְתָה בָזֶה קְדֵשָׁה׃}
}
{
	\eJ{
		And he came back to Judah and said, “I didn’t find her,
		and also the people of the place said There was no sacred
		prostitute here.’”
	}
}
{
	\Table{|p{0.16\linewidth}|p{0.25\linewidth}|p{0.49\linewidth}|}{
		\hline
		\textbf{Proper name} & \textbf{Hebrew} & \textbf{Possible Meaning} \\
		\hline
		Judah & \heb{יְהוּדָה} & See v.1. \\
		\hline
	}
}

\Verse{23}
{
	\hJ{וַיֹּאמֶר יְהוּדָה תִּקַּח לָהּ פֶּן נִהְיֶה לָבוּז הִנֵּה שָׁלַחְתִּי הַגְּדִי הַזֶּה וְאַתָּה לֹא מְצָאתָהּ׃}
}
{
	\eJ{
		And Judah said, “Let her take them or else we’ll be a
		disgrace. Here, I’ve sent this goat kid, and you didn’t
		find her.”
	}
}
{
	\Table{|p{0.16\linewidth}|p{0.25\linewidth}|p{0.49\linewidth}|}{
		\hline
		\textbf{Proper name} & \textbf{Hebrew} & \textbf{Possible Meaning} \\
		\hline
		Judah & \heb{יְהוּדָה} & See v.1. \\
		\hline
	}
}

\Verse{24}
{
	\hJ{וַיְהִי כְּמִשְׁלֹשׁ חֳדָשִׁים וַיֻּגַּד לִיהוּדָה לֵאמֹר זָנְתָה תָּמָר כַּלָּתֶךָ וְגַם הִנֵּה הָרָה לִזְנוּנִים וַיֹּאמֶר יְהוּדָה הוֹצִיאוּהָ וְתִשָּׂרֵף׃}
}
{
	\eJ{
		And it was about three months, and it was told to Judah,
		saying, “Your daughter-in-law Tamar has whored, and, here,
		she’s pregnant by whoring as well.” And Judah said, “Bring
		her out and let her be burned.”
	}
}
{
	\Table{|p{0.16\linewidth}|p{0.25\linewidth}|p{0.49\linewidth}|}{
		\hline
		\textbf{Proper name} & \textbf{Hebrew} & \textbf{Possible Meaning} \\
		\hline
		Judah & \heb{יְהוּדָה} & See v.1. \\
		\hline
		Tamar & \heb{תָּמָר} & See v.6. \\
		\hline
	}
}

\Verse{25}
{
	\hJ{הִוא מוּצֵאת וְהִיא שָׁלְחָה אֶל חָמִיהָ לֵאמֹר לְאִישׁ אֲשֶׁר אֵלֶּה לּוֹ אָנֹכִי הָרָה וַתֹּאמֶר הַכֶּר נָא לְמִי הַחֹתֶמֶת וְהַפְּתִילִים וְהַמַּטֶּה הָאֵלֶּה׃}
}
{
	\eJ{
		She was brought out. And she sent to her father-in-law,
		saying, “I’m pregnant by the man to whom these belong,”
		and she said, “Recognize: to whom do these seal and cords
		and staff belong?”
	}
}
{
%	\footnote{
%		Recognize. Tamar lays down the evidence and uses the same word that
%		Joseph’s brothers used when they laid down the evidence of Joseph’s
%		demise (the bloodstained coat of many colors) in front of Jacob. Like
%		father, like son: Judah might well feel a chill down his back when he
%		hears the word and knows that he was a guilty party on both occasions.
%		We might even imagine that this double sense of his own errors is what
%		moves him to declare: “She’s more right than I am.”
%	}
}

\Verse{26}
{
	\hJ{וַיַּכֵּר יְהוּדָה וַיֹּאמֶר צָדְקָה מִמֶּנִּי כִּי עַל כֵּן לֹא נְתַתִּיהָ לְשֵׁלָה בְנִי וְלֹא יָסַף עוֹד לְדַעְתָּהּ׃}
}
{
	\eJ{
		And Judah recognized and said, “She’s more right than I
		am, because of the fact that I didn’t give her to my son
		Shelah.” And he did not go on to know her again.
	}
}
{
	\Table{|p{0.16\linewidth}|p{0.25\linewidth}|p{0.49\linewidth}|}{
		\hline
		\textbf{Proper name} & \textbf{Hebrew} & \textbf{Possible Meaning} \\
		\hline
		Judah & \heb{יְהוּדָה} & See v.1. \\
		\hline
		Shelah & \heb{שֵׁלָה} & See v.5. \\
		\hline
	}
}

\Verse{27}
{
	\hJ{וַיְהִי בְּעֵת לִדְתָּהּ וְהִנֵּה תְאוֹמִים בְּבִטְנָהּ׃}
}
{
	\eJ{
		And it was at the time that she was giving birth, and here
		were twins in her womb.
	}
}
{
}

\Verse{28}
{
	\hJ{וַיְהִי בְלִדְתָּהּ וַיִּתֶּן יָד וַתִּקַּח הַמְיַלֶּדֶת וַתִּקְשֹׁר עַל יָדוֹ שָׁנִי לֵאמֹר זֶה יָצָא רִאשֹׁנָה׃}
}
{
	\eJ{
		And it was as she was giving birth, and one put out his
		hand, and the midwife took a scarlet thread and tied it on
		his hand, saying, “This one came out first.”
	}
}
{
}

\Verse{29}
{
	\hJ{וַיְהִי כְּמֵשִׁיב יָדוֹ וְהִנֵּה יָצָא אָחִיו וַתֹּאמֶר מַה פָּרַצְתָּ עָלֶיךָ פָּרֶץ וַיִּקְרָא שְׁמוֹ פָּרֶץ׃}
}
{
	\eJ{
		And it was as he pulled his hand back, and here his
		brother came out. And she said, “What a breach you’ve made
		for yourself!” And he called his name Perez.
	}
}
{
	\Table{|p{0.16\linewidth}|p{0.25\linewidth}|p{0.49\linewidth}|}{
		\hline
		\textbf{Proper name} & \textbf{Hebrew} & \textbf{Possible Meaning} \\
		\hline
		Perez & \heb{פָּרֶץ} (\emph{Pāreṣ})
			& “Breach/bursting forth,” from \heb{פרץ}, “break through.” \\
		\hline
	}
}

\Verse{30}
{
	\hJ{וְאַחַר יָצָא אָחִיו אֲשֶׁר עַל יָדוֹ הַשָּׁנִי וַיִּקְרָא שְׁמוֹ זָרַח׃}
}
{
	\eJ{
		And his brother who had the scarlet thread on his hand
		came out after. And he called his name Zerah.
	}
}
{
	\Table{|p{0.16\linewidth}|p{0.25\linewidth}|p{0.49\linewidth}|}{
		\hline
		\textbf{Proper name} & \textbf{Hebrew} & \textbf{Possible Meaning} \\
		\hline
		Zerah & \heb{זָרַח} (\emph{Zāraḥ})
			& “Rising/shining,” from \heb{זרח}, especially of the sun. \\
		\hline
	}
}
